\chapter{Conclusiones y Perspectivas}
\label{ch:conclusion}

\providecommand{\WL}{\textsc{wl}}

Este libro ha construido un \'unico arco. Comenz\'o con la observaci\'on de que
las arquitecturas que hay detr\'as de la revoluci\'on del aprendizaje profundo
---redes completamente conectadas, convolucionales, recurrentes y de
atenci\'on--- fueron concebidas para datos dispuestos sobre dominios euclidianos
regulares, y que esa suposici\'on se rompe ante las mol\'eculas, las redes y las
superficies que impregnan las ciencias. A partir de ah\'i desarroll\'o el
lenguaje matem\'atico de la simetr\'ia (\Cref{ch:foundations}), convirti\'o ese
lenguaje en un principio constructivo de dise\~no (\Cref{ch:priors}), y sigui\'o
el principio a trav\'es de los cinco temas recurrentes ---rejillas, grupos,
grafos, geod\'esicas y gauges
(\Cref{ch:grids,ch:groups,ch:graphs,ch:geodesics,ch:gauges})--- antes de revisar
lo que los modelos resultantes logran en la pr\'actica (\Cref{ch:applications}).

Este cap\'itulo de cierre no introduce maquinaria nueva. Su prop\'osito es atar
los hilos: mostrar que los cap\'itulos describen una sola idea vista desde varios
\'angulos, ser honesto sobre los l\'imites actuales de esa idea y se\~nalar las
direcciones en las que es probable que crezca.

\section{Una visi\'on unificada}
\label{sec:concl:unified}

La lecci\'on m\'as profunda del aprendizaje profundo geom\'etrico es que
arquitecturas que parecen no estar relacionadas son casos particulares de una
sola construcci\'on. Una capa completamente conectada, una convoluci\'on, una red
de grafos y un bloque de autoatenci\'on no difieren en su naturaleza, sino en tres
elecciones geom\'etricas: el \emph{dominio} sobre el que vive la se\~nal, el
\emph{grupo de simetr\'ia} que la arquitectura debe respetar y la
\emph{conectividad} a trav\'es de la cual se permite que fluya la informaci\'on.

El hilo que concreta esta idea es el \emph{paso de mensajes}. Una capa de paso de
mensajes opera sobre un grafo calculando, para cada nodo, un mensaje desde cada
uno de sus vecinos, agregando esos mensajes y actualizando el estado del nodo. El
hecho notable, establecido a lo largo de los cap\'itulos anteriores, es que las
principales arquitecturas del aprendizaje profundo son esta misma capa con grafos
distintos y reglas distintas de compartici\'on de pesos.

\begin{itemize}
  \item Una \emph{capa completamente conectada} es paso de mensajes sobre el grafo
  completo con un peso distinto para cada par de caracter\'isticas. No supone
  estructura geom\'etrica alguna, y por eso mismo carece de sesgo inductivo y debe
  aprender toda invarianza a partir de los datos.
  \item Una \emph{capa convolucional} (\Cref{ch:grids}) es paso de mensajes sobre
  una rejilla regular en la que el peso depende \'unicamente de la posici\'on
  \emph{relativa} del emisor y el receptor. Esa compartici\'on de pesos \emph{es}
  la equivarianza traslacional, y es lo que reduce el n\'umero de par\'ametros de
  cuadr\'atico en el tama\~no de la se\~nal a lineal en el soporte del n\'ucleo.
  \item Una \emph{capa de autoatenci\'on} (el \emph{transformer}) es paso de
  mensajes sobre el grafo completo en el que los pesos se calculan din\'amicamente
  a partir del contenido de los nodos. El patr\'on de comunicaci\'on se aprende
  para cada entrada en lugar de fijarse, lo que compra conectividad global y
  adaptativa a un coste cuadr\'atico en la longitud de la secuencia.
  \item Una \emph{red neuronal de grafos} (\Cref{ch:graphs}) es el caso general:
  paso de mensajes sobre un grafo arbitrario con simetr\'ia de permutaci\'on, del
  cual las tres arquitecturas anteriores son las especializaciones regular,
  completa y adaptativa.
\end{itemize}

Las redes equivariantes a grupos (\Cref{ch:groups}) refinan el caso convolucional
sustituyendo las traslaciones por una simetr\'ia global mayor, de modo que la
compartici\'on de pesos queda gobernada por la acci\'on de un grupo y no por una
rejilla. El aprendizaje en variedades (\Cref{ch:geodesics}) sustituye la rejilla
ambiente por distancias y curvatura intr\'insecas, y recupera el paso de mensajes
como la discretizaci\'on de una difusi\'on continua. Las redes equivariantes gauge
(\Cref{ch:gauges}) reconocen que en un dominio curvo no existe un marco local
can\'onico, y exigen que la construcci\'on sea independiente de esa elecci\'on
arbitraria. En cada paso, el mismo esqueleto ---mensajes locales, pesos
restringidos por la simetr\'ia, composici\'on equivariante--- se lleva a una
geometr\'ia m\'as general.

Esta unificaci\'on no es una mera taxonom\'ia ordenada. Posee aut\'entica fuerza
explicativa y predictiva. Explica \emph{por qu\'e} una arquitectura dada conviene
a un tipo de datos dado: la simetr\'ia del dominio dicta las restricciones de
equivarianza que la red debe satisfacer, y esas restricciones fijan a su vez la
forma de las capas admisibles. Y es predictiva: ante un dominio nuevo, el marco
nos dice qu\'e simetr\'ias deben respetarse y, mediante los resultados de
universalidad equivariante, qu\'e garant\'ias expresivas puede ofrecer el modelo
resultante. El aprendizaje profundo geom\'etrico no es, por tanto, solo una
descripci\'on de las arquitecturas que ya tenemos, sino una receta para dise\~nar
las que a\'un no.

\section{Problemas abiertos}
\label{sec:concl:open}

La misma teor\'ia que unifica estas arquitecturas expone tambi\'en, con inusual
claridad, las fronteras de lo que hoy es posible. Cuatro limitaciones destacan
porque emergen del propio formalismo, y no de una contingencia de ingenier\'ia.

\paragraph{Poder expresivo y la barrera de Weisfeiler--Leman.}
El poder expresivo de las redes est\'andar de paso de mensajes est\'a acotado
exactamente por el test de isomorfismo de grafos de Weisfeiler--Leman
unidimensional: dos grafos que el test no distingue recibir\'an representaciones
id\'enticas de cualquier red de este tipo. Las variantes de orden superior
($k$-\WL{}) rompen este techo, pero su coste crece exponencialmente con $k$, lo
que las hace impracticables incluso para grafos de tama\~no moderado. Dise\~nar
arquitecturas que superen la barrera unidimensional con coste polinomial ---por
ejemplo mediante caracter\'isticas estructurales o atenci\'on de orden superior---
sigue siendo un problema abierto central.

\paragraph{Tensi\'on entre rigidez y adaptabilidad.}
Existe una tensi\'on fundamental entre incorporar la simetr\'ia en la arquitectura
y dejar que se aprenda. Las redes convolucionales equivariantes codifican la
simetr\'ia de forma r\'igida: son eficientes en datos y vienen con garant\'ias
te\'oricas, pero tienen dificultades cuando la simetr\'ia es solo aproximada. Los
mecanismos de atenci\'on son flexibles y aprenden relaciones dependientes del
contexto, pero exigen m\'as datos y ofrecen una regularizaci\'on estructural m\'as
d\'ebil. Hay razones para creer que ambos paradigmas convergen bajo condiciones
apropiadas, pero el punto \'optimo en el espectro rigidez--adaptabilidad para cada
dominio concreto no se comprende todav\'ia.

\paragraph{Equivarianza gauge y topolog\'ia global.}
El formalismo gauge (\Cref{ch:gauges}) da cuenta rigurosa de las simetr\'ias
locales, pero la topolog\'ia global del dominio se le resiste. Los fibrados no
triviales y la holonom\'ia no trivial alrededor de ciclos plantean dificultades a
la vez te\'oricas y computacionales, y las redes equivariantes gauge actuales
dependen de discretizaciones que pueden no preservar las propiedades
topol\'ogicas globales del espacio subyacente. Cerrar la brecha entre la
elegancia del formalismo y las implementaciones eficientes y fieles a la
topolog\'ia constituye un desaf\'io pr\'actico significativo.

\paragraph{Estructuras de orden superior.}
Las herramientas que capturan relaciones m\'as all\'a de las interacciones por
pares ---laplacianos de Hodge, complejos simpliciales y celulares, homolog\'ia
persistente--- son matem\'aticamente maduras, pero las arquitecturas construidas
sobre ellas no lo son. Dise\~nar operaciones de convoluci\'on y atenci\'on sobre
complejos celulares generales que respeten la estructura algebraica de la
topolog\'ia subyacente sin incurrir en costes prohibitivos es un problema abierto
y cada vez m\'as activo.

\section{Direcciones futuras}
\label{sec:concl:future}

Las limitaciones anteriores sugieren, casi por negaci\'on, hacia d\'onde se
dirige el campo. Tres direcciones se desprenden directamente del material de este
libro.

\paragraph{Arquitecturas h\'ibridas convoluci\'on--atenci\'on.}
La correspondencia formal entre convoluci\'on y atenci\'on, junto con el
an\'alisis del espectro rigidez--adaptabilidad, apunta a arquitecturas que
combinen ambas de forma controlada. El objetivo es interpolar entre la eficiencia
en datos de las convoluciones equivariantes y la expresividad de la atenci\'on,
en lugar de elegir una de ellas en bloque. El \'exito del \emph{transformer}
sobre rejillas ---dominios en los que la convoluci\'on ya destaca--- muestra que
la frontera entre ambas es porosa, y una forma principiada de situar un modelo en
cualquier punto de ella ser\'ia de gran valor.

\paragraph{Dominios geom\'etricos complejos.}
El marco de las cinco Ges se extiende de forma natural a espacios producto que
mezclan curvaturas ---hiperb\'olica, esf\'erica y euclidiana--- en un mismo
modelo, lo que conviene a datos simult\'aneamente jer\'arquicos y c\'iclicos. En
la direcci\'on topol\'ogica, la teor\'ia de los haces celulares proporciona una
base para las \emph{sheaf neural networks}, que generalizan las redes de grafos al
dotar a las relaciones entre nodos de la estructura algebraica de un haz. Ambas
extensiones ampl\'ian la clase de dominios que el marco puede abordar conservando
intacto su esqueleto equivariante.

\paragraph{An\'alisis arm\'onico generalizado.}
El an\'alisis arm\'onico ha sido una herramienta unificadora a lo largo del libro,
desde las convoluciones espectrales hasta los arm\'onicos esf\'ericos. Su
resultado central, el teorema de Peter--Weyl, requiere que el grupo de simetr\'ia
sea compacto. Extender el marco a grupos no compactos ---el grupo de Lorentz, los
grupos afines--- exige un an\'alisis arm\'onico no conmutativo que va m\'as all\'a
de lo que proporciona la teor\'ia presente, y ampliar\'ia el alcance de las redes
equivariantes a simetr\'ias relevantes en f\'isica relativista y procesamiento de
se\~nales general.

\section{Reflexi\'on final}
\label{sec:concl:reflection}

El argumento de este libro es que la geometr\'ia no es una colecci\'on de trucos
para exprimir algo m\'as de precisi\'on de un modelo, sino un principio
organizador que explica por qu\'e determinadas arquitecturas funcionan para
determinados datos. Las diferencias entre las redes completamente conectadas, las
convolucionales, las de grafos y los \emph{transformers} reflejan tres elecciones
---dominio, grupo de simetr\'ia y conectividad---, y una vez que esas elecciones
se hacen expl\'icitas, las arquitecturas dejan de ser un zool\'ogico y se
convierten en una familia.

La matem\'atica que sostiene esta visi\'on ---teor\'ia de Lie, geometr\'ia
riemanniana, teor\'ia gauge, topolog\'ia algebraica--- tiene un alcance que va
m\'as all\'a de describir los modelos que ya tenemos. Los resultados de
universalidad equivariante y las correspondencias entre paradigmas dan al marco
una capacidad predictiva: permiten anticipar qu\'e restricciones geom\'etricas
exige un dominio y qu\'e garant\'ias ofrecer\'a la arquitectura resultante. En
ese sentido, el aprendizaje profundo geom\'etrico es menos una teor\'ia acabada
que un programa, y los problemas que este cap\'itulo deja abiertos son una
invitaci\'on a continuarlo. A medida que los datos de los que deseamos aprender
crecen en riqueza geom\'etrica, el argumento de dejar que esa geometr\'ia d\'e
forma al modelo ---y no al rev\'es--- no hace sino fortalecerse~\citep{bronstein2021geometric}.
